\section{Motivation and Objectives}
Performance counters expose properties of execution such as cycles, retired instructions and requests to the memory hierarchy. Their interpretation depends on software context. A function reached through two parent paths can have different data, working sets and reuse behavior. A flat function count merges these paths, while a global counter also includes activity from other tasks and privileged handlers. Hierarchical context records make those distinctions explicit.

\project{} investigates a hardware-assisted connection between normal RISC-V calls and those measurements. The target is an inspectable Rocket implementation integrated with the Chipyard simulation environment~\citep{rocket,chipyard}. Application calls trigger helper execution in the core, rather than requiring an inserted monitoring call at each function. A trusted Linux component controls selection and lifetime, and an offline tool reconstructs a context tree from exported records. This arrangement makes architectural behavior, software ownership and measurement boundaries testable together.

The project has four concrete objectives. First, retain parent/callee identity through nested calls, repeated paths, returns and architectural traps. Second, associate the basic counters and selected native memory-hierarchy events with those contexts while excluding sampled scheduling pauses. Third, implement an executable definition for scalar, stream and irregular memory accesses at retirement. Fourth, provide a complete Linux-to-host result path, including automatic launch of selected ordinary executables and explicit handling of exec, exit and incomplete sessions.

The completed prototype addresses these objectives on one hart. Validation combines directed RTL programs, independently reconstructed retirement observations, actual Linux process execution and host protocol mutations. Performance evaluation uses independently sampled target counters around reproducible bodies. The result is both an implementation and an evidence set: source baselines, ordered patches, configuration hashes, native return codes, shutdown markers and semantic checkers identify what was executed.

The public project entry is \url{https://github.com/Kimberley-Inti/ssa_hpm}. The accepted implementation is tagged \code{v0.2.0}; documentation and this report are a presentation follow-up that preserves its implementation and evidence. Functional navigation groups calling contexts, process accounting, event attribution, export, classification and launch policy, while original development identifiers remain in historical records.

\subsection{Design Basis}
The original project brief describes a procedure-call stack, a parent-sensitive context dictionary and performance records, followed by dataflow objectives for memory patterns and Cache/TLB events~\citep{designbrief}. It proposes a 36\,MiB context region associated with a loader-managed application stack and a small helper-code store named SCCodeTable. Several OS and offline-analysis details are left open.

The implementation makes those open points explicit. Kernel-owned physical pools replace application stack mappings; generated private ROMs implement the helper-code concept; session identity separates tasks and images; and bounded checked streams carry records to host files. Rule 1 defines the address-pattern decision, and the two native profiles define exactly which hardware handshakes increment the fourth counter. These choices form the operating contract described in the following sections.

The accompanying editable report follows the supplied ACL preprint layout and bibliography mechanism~\citep{acltemplate}.
